%----------------------------------------------------------------------------------------
% METADATA BLOCK (YAML)  — see inbox/README.md for every field
%----------------------------------------------------------------------------------------
% uid: ""
% school: "lycee"
% level: "4eme"
% section: "Mathématiques"
% chapter: "Coniques"
% title: "Trois tangentes à une parabole"
% difficulty: "Moyen"
% tags: ["parabole", "tangentes", "orthocentre", "cercle-circonscrit", "directrice"]
% provenance: "adapted"
% batch: "2026-11-20-defis-geometrie-A"
%----------------------------------------------------------------------------------------

% ---STATEMENT---
Le plan est rapporté à un repère orthonormé. On considère la parabole $\mathcal P$ d'équation $x^2=4y$, de foyer $F(0,1)$ et de directrice $\mathcal D:\ y=-1$.
Pour tout réel $u$, on note $T_u$ la tangente à $\mathcal P$ au point d'abscisse $2u$.

Soient $u$, $v$, $w$ trois réels distincts. On note $A=T_v\cap T_w$, $B=T_w\cap T_u$, $C=T_u\cap T_v$ ; ces trois points sont les sommets d'un triangle dont les côtés $(BC)$, $(CA)$, $(AB)$ sont portés par $T_u$, $T_v$, $T_w$. On pose $\sigma_1=u+v+w$, $\sigma_2=uv+vw+wu$, $\sigma_3=uvw$.

% alt: Parabole x²=4y, foyer F, directrice D, trois tangentes formant un triangle ABC, cercle circonscrit passant par F, orthocentre H sur D
\begin{center}
\begin{tikzpicture}[scale=0.9,line cap=round,line join=round,>=stealth]
\draw[->] (-2.8,0) -- (3.9,0) node[right]{$x$};\draw[->] (0,-2.3) -- (0,3.3) node[above]{$y$};
\draw[domain=-2.6:3.7,smooth,samples=60,variable=\t] plot ({\t},{\t*\t/4});
\draw[dashed] (-2.8,-1) -- (3.9,-1) node[right]{$\mathcal D$};
\draw[domain=-2.5:3.6,variable=\t,thin] plot ({\t},{-1*\t-1});
\draw[domain=-2.5:3.6,variable=\t,thin] plot ({\t},{0.5*\t-0.25});
\draw[domain=-2.5:3.6,variable=\t,thin] plot ({\t},{1.6*\t-2.56});
\draw[dotted] (0.95,-0.15) circle (1.492);
\fill(2.1,0.8) circle (1.4pt) node[above right] {$A$};
\fill(0.6,-1.6) circle (1.4pt) node[below right] {$B$};
\fill(-0.5,-0.5) circle (1.4pt) node[left] {$C$};
\fill(0,1) circle (1.4pt) node[above left] {$F$};
\fill(0.3,-1) circle (1.4pt) node[below left] {$H$};
\end{tikzpicture}
\end{center}

\begin{enumerate}
  \item
  \begin{enumerate}
    \item Montrer que $T_u$ a pour équation $y=ux-u^2$.
    \item Montrer que $T_u$ et $T_v$ se coupent au point $(u+v,\,uv)$.
  \end{enumerate}
  \item Montrer que \textbf{l'orthocentre de $ABC$ est le point $H(\sigma_1+\sigma_3,\,-1)$}, qui appartient donc à la directrice $\mathcal D$.
  \item Montrer que \textbf{le cercle circonscrit à $ABC$ passe par le foyer $F$}. (On pourra chercher une équation de ce cercle de la forme $x^2+y^2+Dx+Ey+G=0$ et exprimer $D$, $E$, $G$ à l'aide de $\sigma_1,\sigma_2,\sigma_3$.)
\end{enumerate}

% ---SOLUTION---
\begin{enumerate}[itemsep=14pt]
  \item\leavevmode
  \begin{enumerate}[itemsep=8pt]
    \item $\mathcal P$ est la courbe de $y=\dfrac{x^2}4$ ; $y'=\dfrac x2$.

    \begin{prop}[Tangente]
    La tangente à la courbe de $f$ en $x_0$ a pour équation \ib{$y=f'(x_0)(x-x_0)+f(x_0)$}.
    \end{prop}
    En $x_0=2u$ : $f'(2u)=u$ et $f(2u)=u^2$. Donc $y=u(x-2u)+u^2$, soit \ib{$y=ux-u^2$}.
    \item $T_u\cap T_v$ : $ux-u^2=vx-v^2\iff(u-v)x=u^2-v^2\iff x=u+v$ (car $u\neq v$).

    Alors $y=u(u+v)-u^2=uv$ : le point d'intersection est $(u+v,\,uv)$.
  \end{enumerate}

  \item Les sommets sont $A(v+w,\,vw)$, $B(w+u,\,wu)$, $C(u+v,\,uv)$.

  La droite $(BC)=T_u$ a pour vecteur directeur $\vec d=(1,u)$.

  \begin{prop}[Droite perpendiculaire]
  La droite passant par $A(x_A,y_A)$ et perpendiculaire à un vecteur directeur $(1,u)$ a pour équation \ib{$(x-x_A)+u\,(y-y_A)=0$}.
  \end{prop}
  Hauteur issue de $A$ : $(x-v-w)+u\,(y-vw)=0$.

  On la coupe avec la directrice $y=-1$ : $x=v+w+u(1+vw)=u+v+w+uvw=\sigma_1+\sigma_3$.

  Ce résultat est symétrique en $u,v,w$ : les hauteurs issues de $B$ et de $C$ coupent aussi $\mathcal D$ au même point.

  Les trois hauteurs passent donc par $H(\sigma_1+\sigma_3,-1)$ : \textbf{$H$ est l'orthocentre, et $H\in\mathcal D$}.

  \item \emph{Équation du cercle.} Le cercle $x^2+y^2+Dx+Ey+G=0$ passe par $B(w+u,wu)$ et $C(u+v,uv)$. On soustrait les deux équations :
  \[\big[(u+v)^2-(u+w)^2\big]+\big[(uv)^2-(uw)^2\big]+D(v-w)+E\,u(v-w)=0.\]
  Or $(u+v)^2-(u+w)^2=(v-w)(2u+v+w)$ et $(uv)^2-(uw)^2=u^2(v-w)(v+w)$. On divise par $v-w\neq0$ :
  \[(2u+v+w)+u^2(v+w)+D+Eu=0.\qquad(1)\]
  En échangeant les rôles de $u$ et $v$ (points $A$ et $C$) :
  \[(2v+w+u)+v^2(w+u)+D+Ev=0.\qquad(2)\]
  $(1)-(2)$ : $(u-v)+(u-v)\sigma_2+E(u-v)=0$, car $u^2(v+w)-v^2(w+u)=(u-v)(uv+vw+wu)$. Donc \ib{$E=-(1+\sigma_2)$}.
  Dans $(1)$ : $D=-(2u+v+w)-u^2(v+w)+u(1+\sigma_2)=-\sigma_1+u\big[\sigma_2-u(v+w)\big]=-\sigma_1+u\,vw$, donc \ib{$D=\sigma_3-\sigma_1$}.
  Enfin, avec $C$ : $G=-(u+v)^2-u^2v^2-D(u+v)-Euv$. On remplace $D$ et $E$ et on utilise $\sigma_2=uv+w(u+v)$ :
  \[G=-u^2v^2-uvw(u+v)+w(u+v)+uv+uv\,\sigma_2=w(u+v)+uv=\sigma_2,\]
  puisque $uv\,\sigma_2=u^2v^2+uvw(u+v)$. Ainsi \ib{$G=\sigma_2$}.

  \emph{Le foyer.} Pour $F(0,1)$ :
  \[0^2+1^2+D\times0+E\times1+G=1-(1+\sigma_2)+\sigma_2=0.\]
  $F$ vérifie l'équation du cercle : \textbf{le cercle circonscrit à $ABC$ passe par $F$}. $\blacksquare$
\end{enumerate}

\textbf{Remarque (autre démonstration).} Le projeté orthogonal de $F$ sur $T_u$ est $(u,0)$ : il est sur la tangente au sommet $y=0$. Les trois projetés de $F$ sur les côtés sont donc alignés, et la réciproque de la droite de Simson donne directement $F\in$ cercle circonscrit.

\textbf{Idée :} tout se fait avec le paramètre de tangence : les sommets sont $(u+v,uv)$, les formules sont des polynômes symétriques en $u,v,w$, et l'orthocentre sort sur $y=-1$.

\textbf{Piège :} la tangente en $(2u,u^2)$ a pour pente $u$ (et non $2u$) ; avec le paramètre $2u$ pour l'abscisse, tout est plus propre.
